\begin{table}[!htbp]
\centering
\footnotesize
\caption{Every Signed Bellman Interval and Numerical Account}\label{tab:r16signedfull}
\begin{tabular}{rlrrrrr}
\toprule
$d$ & Statistic & Mean & Lower & Upper & Clip radius & Arithmetic \\
\midrule
10 & $M$ & 0.0012435 & 0.0007427 & 0.0017443 & 0.211739 & 0.00000018 \\
10 & $C$ & -0.0000022 & -0.0006220 & 0.0006175 & 0.266829 & 0.00002456 \\
10 & $G$ & 0.0012457 & 0.0010673 & 0.0014242 & 0.055090 & 0.00002439 \\
50 & $M$ & 0.0013538 & 0.0009967 & 0.0017109 & 0.146951 & 0.00000012 \\
50 & $C$ & 0.0000006 & -0.0005105 & 0.0005117 & 0.214706 & 0.00003080 \\
50 & $G$ & 0.0013532 & 0.0011389 & 0.0015675 & 0.067755 & 0.00003068 \\
\bottomrule
\end{tabular}
\begin{minipage}{0.97\linewidth}\footnotesize
\emph{Notes:} All six registered intervals appear, including inconclusive correction intervals. Endpoints include the additional shared-prefix representation allowance. The centered clipping intervals are fixed before new confirmation and can be asymmetric about zero. The arithmetic column is the interval-sample mean and standard-deviation perturbation allowance, not an empirical population range. Every endpoint is rounded outward.
\end{minipage}
\end{table}
